\documentclass[a4paper,10pt]{article}
\usepackage[utf8]{inputenc}
\usepackage[T1]{fontenc}
\usepackage[french]{babel}
\usepackage[margin=1.5cm]{geometry}
\usepackage{xcolor}
\usepackage{tcolorbox}
\usepackage{enumitem}
\usepackage{multicol}
\usepackage{lmodern}
\usepackage{tabularx}
\usepackage{booktabs}

\renewcommand{\familydefault}{\sfdefault}
\pagestyle{empty}
\setlist{nosep, leftmargin=1.5em}

\definecolor{navy}{RGB}{0,51,102}
\definecolor{teal}{RGB}{0,120,120}
\definecolor{crimson}{RGB}{160,0,0}
\definecolor{green}{RGB}{0,120,50}
\definecolor{lightbg}{RGB}{245,248,252}
\definecolor{orange}{RGB}{200,100,0}
\definecolor{purple}{RGB}{100,0,160}

\tcbset{basebox/.style={
  colback=lightbg, colframe=navy, coltitle=white,
  fonttitle=\bfseries, arc=2mm, boxrule=0.5mm,
  left=3mm, right=3mm, top=2mm, bottom=2mm, after skip=6pt
}}

\newcommand{\key}[1]{\textbf{\textcolor{navy}{#1}}}
\newcommand{\warn}[1]{\textbf{\textcolor{crimson}{#1}}}
\newcommand{\ok}[1]{\textbf{\textcolor{green}{#1}}}

\begin{document}

\begin{center}
  \colorbox{crimson}{\parbox{\dimexpr\textwidth-2\fboxsep}{
    \centering\vspace{3pt}
    {\LARGE\bfseries\textcolor{white}{CYBERSÉCURITÉ — Fiche de révision détaillée}}\\
    {\normalsize\textcolor{white}{BTS SIO — Maël Mignard}}
    \vspace{3pt}
  }}
\end{center}
\vspace{6pt}

\begin{multicols}{2}

% ============================================================
\begin{tcolorbox}[basebox, colframe=crimson, title=1. Triade CIA — Piliers fondamentaux]
\begin{tcolorbox}[colback=red!5, colframe=crimson, boxrule=0.3mm, arc=1mm]
  \textbf{\textcolor{crimson}{C}onfidentialité} : seules les personnes autorisées accèdent aux données.\\
  \textbf{\textcolor{crimson}{I}ntégrité} : les données ne sont pas altérées sans autorisation.\\
  \textbf{\textcolor{crimson}{D}isponibilité} : le service est accessible quand on en a besoin.
\end{tcolorbox}
\vspace{4pt}
\key{Authentification} : vérifier l'identité d'un utilisateur \\
\key{Autorisation} : vérifier ce qu'il a le droit de faire \\
\key{Non-répudiation} : prouver qu'une action a eu lieu
\end{tcolorbox}

% ============================================================
\begin{tcolorbox}[basebox, colframe=crimson, title=2. Menaces et attaques]
\textbf{Attaques sur les données :}
\begin{itemize}
  \item \warn{Phishing} : faux e-mail imitant une entité de confiance $\to$ vol identifiants
  \item \warn{Spear-phishing} : phishing ciblé (sur une personne précise)
  \item \warn{Ransomware} : chiffrement des données + demande de rançon
  \item \warn{Keylogger} : enregistre les frappes clavier
\end{itemize}
\vspace{4pt}
\textbf{Attaques sur les applications web :}
\begin{itemize}
  \item \warn{XSS} (Cross-Site Scripting) : injection JS malveillant dans une page
  \item \warn{SQLi} (Injection SQL) : manipulation de requêtes BDD
  \item \warn{CSRF} : forcer un utilisateur connecté à effectuer une action
  \item \warn{Brute Force} : essai massif de mots de passe
\end{itemize}
\vspace{4pt}
\textbf{Attaques réseau :}
\begin{itemize}
  \item \warn{DDoS} : saturation d'un serveur par des requêtes massives
  \item \warn{Man-in-the-Middle} : interception de communication
  \item \warn{RCE} : Remote Code Execution, exécution distante de code
  \item \warn{Sniffing} : écoute passive du trafic réseau
\end{itemize}
\end{tcolorbox}

% ============================================================
\begin{tcolorbox}[basebox, colframe=navy, title=3. Gestion des accès et identités]
\key{IAM} (Identity and Access Management) :\\
Gérer qui peut accéder à quoi.\\[4pt]
\key{RBAC} (Role-Based Access Control) :\\
Les droits sont liés au rôle, pas à l'individu.\\
Ex: Comptable $\to$ accès finances seulement\\[4pt]
\key{Principe du moindre privilège} :\\
Donner uniquement les droits nécessaires.\\[4pt]
\key{MFA} (Multi-Factor Authentication) :\\
Combiner plusieurs facteurs :
\begin{itemize}
  \item Ce qu'on sait : mot de passe
  \item Ce qu'on a : téléphone (OTP)
  \item Ce qu'on est : empreinte digitale
\end{itemize}
\key{OTP} : One-Time Password (code à usage unique) \\
\key{SSO} : Single Sign-On (une seule authentification) \\
\key{IDP} : Identity Provider (ex: Google, Azure AD)
\end{tcolorbox}

% ============================================================
\begin{tcolorbox}[basebox, colframe=green, title=4. Chiffrement et cryptographie]
\key{Chiffrement symétrique} :\\
Même clé pour chiffrer et déchiffrer.\\
Rapide, mais partage de clé délicat.\\
Ex: AES-256\\[4pt]
\key{Chiffrement asymétrique} :\\
Clé publique (chiffrement) + clé privée (déchiffrement).\\
Ex: RSA, utilisé pour HTTPS, SSH.\\[4pt]
\key{Hashage} : transformation irréversible.\\
Ex: SHA-256, bcrypt (pour les mots de passe).\\
\warn{MD5 et SHA-1 : obsolètes, vulnérables}\\[4pt]
\key{TLS/SSL} : chiffre les communications web \\
\key{HTTPS} = HTTP + TLS \\
\key{PKI} : Infrastructure à clé publique (certificats)\\
\key{Certificat} : lie une identité à une clé publique
\end{tcolorbox}

% ============================================================
\begin{tcolorbox}[basebox, colframe=orange, title=5. RGPD — Règlement Général Protection Données]
\key{En vigueur depuis} : 25 mai 2018 (UE)\\[4pt]
\textbf{Droits des personnes :}
\begin{itemize}
  \item Droit d'accès à ses données
  \item Droit de rectification
  \item Droit à l'effacement (« droit à l'oubli »)
  \item Droit à la portabilité
  \item Droit d'opposition au traitement
\end{itemize}
\vspace{4pt}
\textbf{Obligations des entreprises :}
\begin{itemize}
  \item Consentement explicite et éclairé
  \item Notification de violation de données : \warn{72h max}
  \item Désigner un \key{DPO} si traitement à grande échelle
  \item \textit{Privacy by design} : sécurité intégrée dès la conception
\end{itemize}
\vspace{4pt}
\textbf{Sanctions :}\\
Jusqu'à \warn{4\% du CA mondial} ou 20M€ \\
\key{CNIL} : autorité de contrôle en France
\end{tcolorbox}

% ============================================================
\begin{tcolorbox}[basebox, colframe=purple, title=6. PCA / PRA — Continuité et reprise]
\key{PCA} (Plan de Continuité d'Activité) :\\
Maintenir l'activité malgré un incident majeur.\\[4pt]
\key{PRA} (Plan de Reprise d'Activité) :\\
Remettre en marche le SI après un sinistre.\\[4pt]
\textbf{Indicateurs clés :}
\begin{itemize}
  \item \key{RTO} (Recovery Time Objective) : durée maximale d'interruption tolérée
  \item \key{RPO} (Recovery Point Objective) : perte de données maximale acceptable (en temps)
\end{itemize}
\vspace{4pt}
\textbf{Risques couverts :} incendie, inondation, panne serveur, cyberattaque, erreur humaine.\\[4pt]
\key{Site de secours} : site miroir ou cloud de repli\\
\key{Redondance} : dupliquer les équipements critiques
\end{tcolorbox}

% ============================================================
\begin{tcolorbox}[basebox, colframe=navy, title=7. Analyse des risques]
\key{Risque} = Menace $\times$ Vulnérabilité $\times$ Impact\\[4pt]
\textbf{Méthode EBIOS Risk Manager (ANSSI)} :\\
Identifier, évaluer et traiter les risques.\\[4pt]
\textbf{Niveaux de criticité :}
\begin{tabular}{ll}
\textbf{Probabilité} & Faible / Moyenne / Forte \\
\textbf{Impact} & Mineur / Significatif / Critique \\
\end{tabular}
\vspace{4pt}\\
\textbf{Traitements possibles :}
\begin{itemize}
  \item \ok{Réduction} : mesures de sécurité
  \item Transfert : assurance, externalisation
  \item Acceptation : risque trop faible pour agir
  \item Évitement : arrêter l'activité risquée
\end{itemize}
\end{tcolorbox}

% ============================================================
\begin{tcolorbox}[basebox, colframe=teal, title=8. Sécurité des réseaux]
\key{Firewall} : filtre le trafic entrant/sortant \\
\key{IDS} (Intrusion Detection System) : détecte les intrusions \\
\key{IPS} (Intrusion Prevention System) : bloque les intrusions \\
\key{DMZ} (Demilitarized Zone) : zone réseau semi-sécurisée (serveurs web publics) \\
\key{VPN} : tunnel chiffré sur réseau public \\
\key{VLAN} : isolation logique du réseau \\[4pt]
\textbf{Architecture en couches :}\\
Internet $\to$ Firewall $\to$ DMZ $\to$ Firewall $\to$ LAN interne
\end{tcolorbox}

% ============================================================
\begin{tcolorbox}[basebox, colframe=orange, title=9. Gestion de crise — Réponse à incident]
\textbf{Étapes de réponse à un incident :}
\begin{enumerate}
  \item \key{Détection} : identifier l'incident
  \item \key{Confinement} : isoler les systèmes compromis
  \item \key{Éradication} : supprimer la menace
  \item \key{Restauration} : remettre en service
  \item \key{Retour d'expérience} : documenter, améliorer
\end{enumerate}
\vspace{4pt}
\key{CERT/CSIRT} : équipe de réponse à incident \\
\key{Forensique} : analyse post-incident des preuves \\
\warn{Toujours documenter les actions menées pendant la crise}
\end{tcolorbox}

% ============================================================
\begin{tcolorbox}[basebox, colframe=green, title=10. Bonnes pratiques de sécurité]
\textbf{Pour les utilisateurs :}
\begin{itemize}
  \item Mot de passe long et unique (gestionnaire de mots de passe)
  \item Activer le MFA sur tous les comptes sensibles
  \item Ne jamais cliquer sur un lien suspect
  \item Mettre à jour régulièrement les logiciels (patchs)
\end{itemize}
\vspace{4pt}
\textbf{Pour les développeurs :}
\begin{itemize}
  \item Valider et échapper toutes les entrées utilisateurs
  \item Utiliser HTTPS partout
  \item Ne jamais stocker les mots de passe en clair (utiliser bcrypt)
  \item Principe du moindre privilège pour les comptes BDD
\end{itemize}
\vspace{4pt}
\textbf{Pour les administrateurs :}
\begin{itemize}
  \item Journalisation et surveillance des logs
  \item Segmentation réseau (VLAN, DMZ)
  \item Tests de pénétration réguliers (pentest)
  \item Politique de sauvegarde 3-2-1
\end{itemize}
\end{tcolorbox}

\end{multicols}

\vspace{4pt}
\begin{tcolorbox}[basebox, colframe=crimson, title=Points clés à retenir pour l'examen]
\begin{multicols}{2}
\begin{itemize}
  \item \textbf{CIA} : Confidentialité, Intégrité, Disponibilité
  \item \textbf{RBAC} = droits liés au rôle, pas à l'individu
  \item \textbf{MFA} = ce qu'on sait + ce qu'on a + ce qu'on est
  \item \textbf{RGPD} : notification violation = \warn{72 heures}
  \item \textbf{PCA} = maintenir ; \textbf{PRA} = reprendre après sinistre
  \item \textbf{RTO} = durée max d'interruption ; \textbf{RPO} = perte max de données
  \item \textbf{XSS} = injection JS ; \textbf{SQLi} = injection SQL
  \item \textbf{Phishing} = usurpation d'identité par mail
  \item \textbf{Chiffrement asymétrique} : clé publique/privée (RSA)
  \item \textbf{Hashage} = irréversible (SHA-256, bcrypt)
  \item \textbf{DMZ} = zone exposée internet, entre 2 firewalls
  \item \textbf{DPO} = responsable RGPD dans l'entreprise
\end{itemize}
\end{multicols}
\end{tcolorbox}

\end{document}
